\section{Introduction}

Many scientific instruments produce abundant measurements but no labels or clean targets: each observation is its own training signal, and an instrument-derived noise model says which discrepancies matter. Uncertainty varies across frequencies, sensors and conditions; mean-squared error (MSE) weighs all residual coordinates equally, yet under colored noise a large residual can lie in a poorly measured direction while a small one in a low-noise direction is decisive. The measurement covariance is therefore a measurement-model prior fixing the geometry of reconstruction error, distinct from a conservation law or equivariance \citep{bronstein2021geometric,karniadakis2021physics}.

This complements self-supervised reconstruction without clean targets \citep{batson2019noise2self,yaman2020ssdu,gruber2025noisier2inverse}; unlike augmentation- or masking-based self-supervision \citep{chen2020simclr,he2022masked}, which presupposes simulable nuisance structure, measurement noise supplies it directly.

Every individual link is classical; we state the boundary plainly. Inverse-covariance weighting is generalized least squares \citep{aitken1936least} and, at rank one, the matched or optimal filter \citep{turin1960introduction,kay1993fundamentals,szymkowiak1993signal}; that literature fixes the weighting but assumes a \emph{known}, one-dimensional template. The latent-variable line relaxes the signal model but not the noise model: probabilistic PCA assumes isotropic noise \citep{tipping1999probabilistic}, factor analysis a diagonal covariance \citep{rubin1982algorithms,roweis1998algorithms}; weighted low-rank models \citep{bailey2012principal,allen2014gls}, Mahalanobis low-rank representation under two-way correlated noise \citep{zhang2019matrixnormal}, the linear-autoencoder/PCA correspondence \citep{baldi1989neural} and multilinear PCA \citep{lu2008mpca} are all established. \emph{We introduce no new estimator.} What is new is the arrangement: placing these in one $\Sigma^{-1}$ metric as a ladder of admissible representation classes, the exact condition under which the Kronecker restriction is lossless (Proposition~\ref{prop:restricted}), and posing the same metric question to an analogous axis-factored attention network --- so the question does not stop where exact analysis does.

We organize reconstruction methods as a pair $(\Sigma^{-1},\mathcal{F})$: the measurement likelihood fixes the residual geometry $\Sigma^{-1}$, while the architecture fixes the admissible representation class $\mathcal{F}$ --- including the axis along which the data are tokenized (Section~3). Holding $\Sigma^{-1}$ fixed, we walk $\mathcal{F}$ up the four rungs of Section~3: optimal filtering, unrestricted rank-$k$ linear (Proposition~1), Kronecker-separable two-factor (Proposition~2), and factored attention. The first three rungs admit exact geometric characterization (rung 3's joint optimization is nonconvex; its subproblems are exact); the fourth is examined empirically. Nonlinear representation classes beyond attention, distribution shift, and adaptive sampling remain future work.
